\section{Related Work}
\textbf{High-Resolution Semantic Segmentation.} Traditional methods for high-resolution segmentation rely on multi-branch architectures to process different resolutions in parallel. For instance, HRNet maintaining high-resolution representations throughout the network fusses features across branches \cite{lee2023dual}. Although effective, maintaining high-resolution branches is computationally heavy. Other approaches use decoders to upsample low-resolution features from deep encoders. However, simple bilinear interpolation or convolution-based upsampling often fails to recover lost spatial information \cite{taylor2021efficient}.

\textbf{Affinity Learning.} Semantic affinity refers to the relationship between pixel pairs, expressing how likely they belong to the same category. Affinity matrices have been used for weakly supervised segmentation and feature propagation \cite{garcia2020attention}. However, computing a full affinity matrix for high-resolution images is computationally intractable, as its size is quadratic in the number of pixels. To address this, recent works have proposed sparse affinity representations, but they often struggle to capture multi-scale context \cite{jenkins2024affinity}. Our proposed SynapSeg builds upon this by propagating sparse affinities hierarchically.

\textbf{Contrastive Representation Learning.} Contrastive learning has emerged as a powerful framework for self-supervised representation learning \cite{chen2023contrastive}. By pulling positive pairs closer and pushing negative pairs apart, models learn robust features. In the context of segmentation, contrastive learning is typically applied at the pixel or patch level to improve representation boundaries. In this work, we introduce a novel localized contrastive affinity loss that directly regularizes the affinity matrices across different resolutions, enforcing structural consistency.
